% Ad Quintum Fratrem 2.6 --- headnote
% ad-quintum-fratrem-02-06, mid-May 56 BC, written at Rome

\textit{Marcus to Quintus, written at Rome a few days after the
Ides of May 56 BC. Quintus is in Sardinia on Pompey's grain
commission; Cicero is back from the Antium retreat where the
\textit{Att.} 4.4A--4.8 cluster was written, and the political
weather has shifted in his favour twice over.}

\textit{The first piece of news is the Senate's refusal of a
\textit{supplicatio} for the consul of 58 --- Aulus Gabinius,
Cicero's old enemy --- on his Judaean campaign as governor of
Syria. The denial of a thanksgiving to a returning general was
all but unprecedented (\textit{Procilius adiurat hoc nemini
accidisse}), and applauded outside the curia. Cicero's relish
is sharpened by his being away when the vote came: ``it is a
clean verdict \textit{[Greek: eilikrines]}, without my fighting
and without my favour.''}

\textit{The second piece is the matter that lies under his
silence: the \textit{ager Campanus}, the Campanian public land
distributed under Caesar's consular law of 59. The agenda for
the Ides and the day after had marked it for debate; the debate
did not happen. Cicero's own line is the famous idiom
\textit{in hac causa mihi aqua haeret} --- the clepsydra of the
courts has stopped: in this case my water sticks. After Luca
he cannot speak against the distribution; the older commitments
do not let him speak for it. The rest is private and waits for
the brother's return: come quickly, the boys ask, and you will
dine with us when you come.}
